% ==============================
% Executive Summary
% ==============================

\chapter*{Executive Summary}
\addcontentsline{toc}{chapter}{Executive Summary}

Quantum computing today suffers from an acute hardware fragmentation problem. Commercial processors built by IBM, Google, Rigetti, and IonQ rely on fundamentally different physical technologies, ranging from superconducting transmon circuits and trapped ions to neutral atoms. Each architecture imposes its own native gate basis, physical coupling constraints, and decoherence profiles. Consequently, a quantum circuit designed and optimized for an IBM heavy-hex processor cannot execute on a linear ion-trap machine without substantial manual refactoring. Developers are forced to act as human compilers, inserting SWAP operations by hand to bridge non-adjacent qubits. This tight coupling between algorithm design and target hardware creates severe vendor lock-in and raises the entry barrier for students and researchers.

This Final Year Project presents the \textbf{Quantum Virtual Machine (QVM)}, an open-source, hardware-agnostic execution runtime that brings the ``Write Once, Run Anywhere'' (WORA) paradigm to quantum computing. Much like the Java Virtual Machine decoupled bytecode from physical microprocessors, QVM decouples quantum algorithm logic from physical chip topologies. The system ingests standard quantum circuit descriptions (OpenQASM 2.0, OpenQASM 3.0, and JSON format), validates their grammar using an AST-based parser, decomposes non-native operations into target primitives, and automatically routes two-qubit interactions across restricted coupling graphs.

At the core of the transpilation engine, QVM offers two distinct routing strategies: a fast Greedy Breadth-First Search (BFS) router for rapid prototyping, and a lookahead SABRE-inspired router (SWAP-based Bidirectional Heuristic Search) that evaluates decay-weighted distances across future gate layers. In benchmarks, the SABRE heuristic reduces overall SWAP overhead and circuit depth by 30\% to 40\% compared to naive greedy routing, preserving circuit fidelity on noise-constrained backends.

To validate circuit behavior without requiring paid cloud access or waiting in multi-hour hardware queues, QVM integrates dual simulation engines:
\begin{enumerate}
    \item \textbf{Dense Statevector Simulator:} Executes exact unitary evolution for circuits up to 12 qubits using vectorized NumPy array slicing and multi-dimensional tensor contractions. It models mid-circuit measurements with statevector collapse and maintains classical memory registers for hybrid feed-forward logic.
    \item \textbf{Matrix Product State (MPS) Simulator:} Simulates larger quantum systems scaling to 20+ qubits for low-entanglement circuits by representing states as one-dimensional tensor trains and truncating singular values via Singular Value Decomposition (SVD).
\end{enumerate}

To bridge the gap between idealized simulation and physical hardware, the runtime also includes an open-system Kraus noise simulator. It models depolarizing channels, energy decay ($T_1$ relaxation), pure dephasing ($T_2$ decoherence), and asymmetric readout errors, providing pre-calibrated emulations of 5-qubit and 7-qubit devices.

Users interact with QVM through three distinct channels: a command-line interface (CLI) for shell scripting and local file compilation; a FastAPI-based REST service with asynchronous execution endpoints and Supabase cloud persistence; and an interactive web dashboard developed with Next.js and React that renders live SVG circuit diagrams, probability distributions, and KaTeX-formatted quantum amplitudes.

The system was implemented in Python 3.10+ and validated against 215 automated tests across 24 test suites, covering unit operations, gate decompositions, routing invariants, and end-to-end execution of canonical quantum algorithms (Bell state generation, GHZ entanglement, Bernstein-Vazirani parity search, and Grover's unstructured search). Running on standard consumer laptops, QVM executes 10-qubit circuits in under one second while maintaining an installation footprint below 50~MB. By providing a clean, transparent, and modular execution environment, QVM equips students and researchers with a practical platform for exploring quantum compiler mechanics, algorithm design, and hardware-agnostic computing.

\vspace{0.8cm}

\noindent\textbf{Keywords:} Quantum Virtual Machine, Hardware Abstraction, Transpilation, Qubit Routing, SABRE Heuristic, Statevector Simulation, Matrix Product States, OpenQASM 3.0, Noise Modeling.

\vspace{0.5cm}

\noindent\textbf{Key System Metrics:}
\begin{itemize}[leftmargin=*, noitemsep]
    \item \textbf{Implementation Language:} Python 3.10+ (NumPy, Lark, FastAPI, Matplotlib)
    \item \textbf{Frontend Stack:} Next.js 16, React 19, TypeScript, Tailwind CSS
    \item \textbf{Supported Formats:} OpenQASM 2.0, OpenQASM 3.0 (with loops and classical conditions), JSON
    \item \textbf{Routing Strategies:} Greedy BFS, SABRE lookahead routing with decay weighting
    \item \textbf{Simulation Backends:} Exact Statevector (up to 12 qubits), Matrix Product State (20+ qubits)
    \item \textbf{Noise Channels:} Depolarizing, Amplitude Damping ($T_1$), Phase Damping ($T_2$), Readout confusion matrix
    \item \textbf{Test Suite:} 215 automated pytest tests passing across 24 modules
    \item \textbf{Algorithms Verified:} Bell State ($|\Phi^+\rangle$), GHZ-3, Bernstein-Vazirani, Grover Search ($N=8$)
\end{itemize}
